% ===== 原文 pp.32–34 =====
% [原文 p.32]

这里 $m_i$ 是 $\SU(2)_i$ 的质量参数，$Z^{\text{instanton}}_{\text{n.p.}}(a,m_i,q)$ 是来自北极的
瞬子贡献，$Z^{\text{gauge,one-loop}}_{S^4}(a)$ 是规范多重态给出的单圈贡献。先从规范单圈因子
开始，它是
\begin{equation}
Z^{\text{gauge,one-loop}}_{S^4}(a)=\prod_{\pm}\frac{1}{\Gamma_2(\pm a|\epsilon_1,\epsilon_2)
\Gamma_2(\pm a+\epsilon_1\epsilon_2|\epsilon_1,\epsilon_2)}=\Upsilon(a)\Upsilon(Q-a).
\tag{5.19}
\end{equation}

% [Tutor Note: p.32, (5.19) 分母中第二个 $\Gamma_2$ 的辐角原文印作 $\pm a+\epsilon_1\epsilon_2$，
%  这里 $\epsilon_1\epsilon_2$ 是乘积。把它与 §4.3.2 的式 (4.18) 对照：
%  $Z^{G}_{\text{one-loop}}=\prod_{\alpha>0}1/[\Gamma_2(\alpha\cdot a+\epsilon_1+\epsilon_2|
%  \epsilon_1,\epsilon_2)\Gamma_2(\alpha\cdot a|\epsilon_1,\epsilon_2)]$，同是规范多重态的单圈
%  因子，那里的第二项是 $\epsilon_1+\epsilon_2$。故 (5.19) 的 $\epsilon_1\epsilon_2$ 应为
%  $\epsilon_1+\epsilon_2$（已用 800 dpi 裁图复核，原文确实印的是 $\epsilon_1\epsilon_2$）。
%  正文照原文排印。]

连同因子 $Z^{\text{trifund}}_{S^4}(m_1,m_2,a)$ 与 $Z^{\text{trifund}}_{S^4}(a,m_3,m_4)$，这就给出
了 Liouville 四点函数 (2.59) 表达式中三点函数的乘积。

这强烈地提示其余的因子也应当对得上，也就是说：(2.59) 中的共形块 $F$ 与瞬子配分函数
$Z^{\text{instanton}}_{\text{n.p.}}$ 应当相等：
\begin{equation}
F_{\alpha_1,\alpha_2,ip+Q/2,\alpha_3,\alpha_4}(q)\quad\text{``$=$''}\quad
Z^{\text{instanton}}_{\text{n.p.}}(a,m_i,q)
\tag{5.20}
\end{equation}
在如下识别下
\begin{equation}
\alpha_i=m_i+\frac{Q}{2},\qquad ip+\frac{Q}{2}=a+\frac{Q}{2}.
\tag{5.21}
\end{equation}

等式 (5.20) 在如下意义下「几乎」成立。正如 4.3 节已经提到的，我们只能计算 $\U(2)$ 规范场的
瞬子贡献；对 $\U(2)$ 而言并不存在半超多重态，因为它的双重态与反双重态并不相同。三基表示
$\SU(2)_1\times\SU(2)_2\times\SU(2)$ 中质量参数为 $m_1,m_2$ 的半超多重态，可以看作 $\U(2)$
中两个基础超多重态，其质量分别为 $m_1+m_2$ 与 $-m_1+m_2$。这两个基础超多重态的贡献可以用
$\U(2)\times\U(2)$ 的双基表示的公式 (4.24) 算出。类似地，我们把三基表示
$\SU(2)\times\SU(2)_3\times\SU(2)_4$ 中的半超多重态看作 $\U(2)$ 中两个基础超多重态，质量
分别为 $m_3+m_4$ 与 $m_3-m_4$。

于是 (5.20) 的具体形式为
\begin{equation}
Z^{\text{instanton}}_{\text{n.p.},\ \U(2)\ \text{formulation}}(a,m_i,q)
=(1-q)^{2(m_2-Q/2)(Q/2-m_3)}F_{\alpha_1,\alpha_2,ip+Q/2,\alpha_3,\alpha_4}(z).
\tag{5.22}
\end{equation}

其中 $(1-q)$ 的分数幂因子，就是 4.3 节已经提到的那个伪贡献。这里建议读者把上式的两边算到
$q^2$ 阶来核对等式。这只需用到本综述中已经列出的公式即可完成。\footnote{作者较晚才加入
导致文献 [2] 的那项合作，其贡献只是提供了进行瞬子计数的 Mathematica 代码，用来核对这一
等式。}

对物理上正确的配分函数版本，应当把这个伪预因子去掉，于是我们直接令
\begin{equation}
Z^{\text{instanton}}_{\text{n.p.},\text{corrected}}(a,m_i,q)
=F_{\alpha_1,\alpha_2,ip+Q/2,\alpha_3,\alpha_4}(z).
\tag{5.23}
\end{equation}
于是我们有等式
\begin{equation}
Z_{S^4}(m_1,m_2;m_3,m_4;\tau)=\langle V_{\alpha_1}(\infty)V_{\alpha_2}(1)V_{\alpha_3}(z)
V_{\alpha_4}(0)\rangle_{\Liouville}
\tag{5.24}
\end{equation}

% [原文 p.33]

其中等号又一次成立到乘以一个与 $m_i$ 和 $z$ 都无关的函数为止。\footnote{例如 (5.18) 的被积
函数含 $|q|^{2a^2}$，而 (2.59) 的被积函数含 $|q|^{2a^2+Q^2/2}$。因此粗略看来，两个表达式
之间相差（除其它因子外）一个 $|q|^{Q^2/2}$。但要注意，$S^4$ 这类弯曲空间上的 $\cN=2$
超对称理论具有局部超对称抵消项 [32--34]，它们可以把配分函数乘以形如 $|f(\tau)|$ 的因子，
其中 $f$ 是某个全纯函数。正是这一不确定性解释了 $|q|^{Q^2/2}$ 的差异。}

% [Tutor Note: p.32, (5.22) 里 $(1-q)^{2(m_2-Q/2)(Q/2-m_3)}$ 这种分数幂正是 §4.3.3（原文
%  p.26）所说的「U(1) 因子」残余：用 $\U(2)$ 而非 $\SU(2)$ 作瞬子计数时，双重态与反双重态
%  不同，于是多出一个没有物理意义的整体因子，(5.23) 把它去掉后才与 (2.59) 的共形块对应。
%  另请留意：本译本的脚注由 LaTeX 自动编号，与原文字号（原文此页为脚注 3、下一页为脚注 4）
%  不必一一对应，回查原文时请以「页码 + 内容」为准。]

% [Tutor Note: p.33, 脚注 4 的指数差值得补一句推导：由 (2.55) 之后紧接的定义给出的
%  $\Delta_p=p^2+Q^2/4$ 可知 (2.59) 的被积函数带 $|q|^{2\Delta_p}=|q|^{2p^2+Q^2/2}$，而 (5.18)
%  带的是 $|q|^{2a^2}$，两者之差中与 $Q$ 有关的那一项正是 $|q|^{Q^2/2}$。这类只依赖 $\tau$
%  的整体因子正是 (5.17)、(5.24) 里「等号只到与 $m_i$ 无关的函数为止」的具体含义，因为
%  $\tau$ 与 $m_i$ 无关。]

这一分析可以推广到规范群为 $\SU(2)^n$ 的线性箭图（quiver）与环形箭图。限制来自我们无力计算这样
一种瞬子配分函数：它含有一个三基表示半超多重态，而该三基表示的三个 $\SU(2)$ 全都耦合到动力学
规范场。

\subsection{六维解释}

我们把目前看到的东西小结一下。第 3 节中，我们引入了一类与带穿刺点的 Riemann 曲面相联系
的四维 $\cN=2$ 理论。有 $n$ 个穿刺点时，相应的理论具有 $\SU(2)^n$ 味对称性。第 5.1 节中，
我们借助第 4.2 节回顾的公式，计算了这些理论在 $S^1\times S^3$ 上的配分函数。我们发现，它
等于我们在第 2.1 节回顾过的、定义在同一 Riemann 曲面上的二维 $q$-形变 $\SU(2)$ 杨--米尔斯
理论。

第 5.2 节中，我们借助第 4.3 节回顾的公式，计算了这些理论在 $S^4$ 上的配分函数。我们发现，
它等于我们在第 2.2 节回顾过的、定义在同一 Riemann 曲面上的 Liouville 理论。我们应当如何
理解这些对应关系？第 1 节已经给出了梗概。这里我们稍微多看一些细节。

\subsubsection{六维 $\cN=(2,0)$ 理论与四维 class S 理论}

我们从 $\SU(2)$ 型的六维 $\cN=(2,0)$ 理论出发，把它放在一个 Riemann 曲面 $C$ 上。先讨论
这个理论对我们重要的几个性质。当它在 $S^1$ 上紧致化时，就变成五维的最大超对称 $\SU(2)$
杨--米尔斯理论。它究竟在什么确切意义上「变成」五维理论，目前仍有激烈争论。一个奇怪的事实
是：五维理论的瞬子数，应当与六维理论在 $S^1$ 上的 Kaluza--Klein 动量等同。实际使用中已知：
就超对称所保护的那些量而言，我们只需在计算中计入超对称的瞬子贡献即可。关于这一点的更多
细节见本文集中的 [35] 章。我们绝对不应把五维的 Kaluza--Klein 有质量场塔也一并计入，否则
会导致重复计数。

% [Tutor Note: p.33, 「Contribution [35]」指本文集（Localization techniques in quantum field
%  theories，共 17 章，见 [1]）中收录的另一章，而不是一篇独立论文。原文这段的要点是「该计入
%  什么」：既然上一句已经把五维瞬子数与六维在 $S^1$ 上的 KK 动量等同，那么五维的
%  Kaluza--Klein 有质量塔就绝不能再加进来——那正是同一批态被数了两遍，即「overcounting」。]

这个六维理论有余维 2 的算符：在每一个这样的算符上，我们都带有一个 $\SU(2)$ 味对称性。读者
想必早已熟悉支在点上的算符：计算这类点算符的关联函数，几乎是我们学习量子场论时最先做的事。
支在线上的算符也相当熟悉：在规范群为 $G$ 的四维规范理论中，对 $G$ 的任意表示 $R$，我们都
可以取规范场沿一条线 $L$ 的路径排序指数的迹。

% [原文 p.34]

这就定义了 $L$ 上表示 $R$ 的 Wilson 线。它引入了世界线为 $L$ 的外加带电粒子。类似地，我们
可以在给定的世界线上引入外加的带磁荷粒子：这就确定了 't Hooft 圈算符。六维理论的余维 2
算符与之类似：它沿六维时空中的四个方向延展。这个四维子空间可以带有自己的外 $\SU(2)$ 背景
场，与其味对称性相耦合。

于是，给定一个赋有度规的 Riemann 曲面 $C$ 以及其上 $n$ 个选定的点 $p_i$，我们考虑形如
$\R^{1,3}\times C$ 的时空以及其上的 $\SU(2)$ 型 $\cN=(2,0)$ 理论，并对每个 $i$ 在
$\R^{1,3}\times p_i$ 上放一个余维 2 算符。把 $C$ 的面积取得很小，这就定义了味对称性为
$\SU(2)^n$ 的四维理论。为了保留一定数量的超对称，我们利用 $\cN=(2,0)$ 理论的
$\SO(5)_R$ 对称性。$C$ 的曲率使 $C$ 上的旋量丛变得非平凡；但六维理论的超荷不仅在 $C$ 的
$\SO(2)$ 度规曲率下带荷，也在 $\SO(5)_R$ 下带荷。于是我们取 $\SO(5)_R$ 的一个 $\SO(2)$
子群，并在那里放置补偿曲率，使得有一部分超荷成为 $C$ 上平凡丛的截面。这种方法通常称为
部分拓扑扭转。

% [Tutor Note: p.34, 部分拓扑扭转的计数原文没有给出：$\cN=(2,0)$ 理论有 16 个超荷；$C$ 的
%  度规曲率与 $\SO(2)_R$ 的补偿曲率方向相反、相互抵消后，只有一半（8 个）超荷成为 $C$ 上
%  平凡丛的截面，而这 8 个正是四维 $\cN=2$ 的超荷数。所谓「部分」，就是只用了
%  $\SO(5)_R$ 的一个 $\SO(2)$ 子群，而不是整个 $\SO(5)_R$。]

最朴素的选择是取子群 $\SO(2)_R\times\SO(3)_R\subset\SO(5)_R$，并用其中的 $\SO(2)_R$ 来做
部分拓扑扭转。所得理论在四维中是 $\cN=2$ 超对称的。这种部分拓扑扭转还有一个特点：大部分
可计算的四维物理只依赖于总面积以及该二维曲面的复结构，而不依赖于度规的具体选取。

这使我们能进行如下操作：把带穿刺的 Riemann 曲面分解成由细长圆柱连接起来的三穿刺球面。
对每一根圆柱，我们先绕其 $S^1$ 作约化。这样剩下的就是一条边上的五维 $\SU(2)$ 超杨--米尔斯
理论，它给出一个四维 $\SU(2)$ 规范多重态。可以论证，所得规范多重态的复化耦合由 $\tau$ 给出，
即几何上与圆柱相联系的那个复参数。对每个三穿刺球面，我们得到某个具有 $\SU(2)^3$ 对称性的
四维 $\cN=2$ 理论：这就是我们最青睐的那个理论，即 $\SU(2)^3$ 的三基表示中的半超多重态。

如此得到的四维理论 $\cS(C)$，由带穿刺的二维 Riemann 曲面 $C$ 内在地决定。然而，
Riemann 曲面有多种分解成三穿刺球面的方式。每一种这样的分解都给出一种四维拉格朗日描述。
从亏格为 $g$、有 $n$ 个穿刺点的曲面出发，我们总有 $3g-3+n$ 个 $\SU(2)$ 规范多重态和
$2g-2+n$ 个三基表示中的半超多重态，但两种不同分解的 $3g-3+n$ 个耦合常数以复杂的方式相
关联，而各个三基表示的 $n$ 个味对称性也以一种有趣的方式被置换。最常见的情形是：一种拉格朗日量
描述中的基本场，在另一种拉格朗日描述中表现为磁单极或磁单极束缚态。

% [Tutor Note: p.34, $3g-3+n$ 与 $2g-2+n$ 的来历（原文省略）：把 $C$ 分解成三穿刺球面后，
%  每个「顶点」（三穿刺球面）给一个三基表示半超多重态，每条「内边」（细管）给一个 $\SU(2)$
%  矢量多重态；对带 $n$ 条外腿的三价图用 Euler 公式即得顶点数 $2g-2+n$、内边数 $3g-3+n$。
%  自检：四穿刺球面（$g=0$，$n=4$）给 2 个三穿刺球面、1 根管；六穿刺球面（$g=0$，$n=6$）给
%  4 个三穿刺球面、3 根管，正与图 12 的「四个三穿刺球面 + 三根管子」以及下文「规范群
%  $\SU(2)^3$、四个三基表示」一致。]

作为示例，考虑一个有六个穿刺点的球面，见图 12。它可以由四个三穿刺球面通过三根管子连接
而成。因此规范群是 $\SU(2)^3$，并且有四个三基表示。
